\documentclass[10pt,a4paper]{amsart}
\usepackage[margin=19mm]{geometry}
\usepackage{amsmath,amssymb,mathtools}
\usepackage[scr=rsfs,cal=euler]{mathalfa}
\usepackage{xfrac}
\usepackage[unicode,hidelinks,bookmarks=false]{hyperref}
\newcommand{\ie}{i.e.\ }
\newcommand{\cf}{cf.\ }
\newcommand{\resp}{respectively\ }
\newcommand{\Subsection}{\subsection}
\providecommand{\nobreakdash}{\nobreak\hbox{-}}
\setlength{\emergencystretch}{2em}
\multlinegap0pt
\usepackage{booktabs}
\newtheorem{theorem}{Theorem}[section]
\newtheorem{lemma}[theorem]{Lemma}
\newtheorem{proposition}[theorem]{Proposition}
\newtheorem{corollary}[theorem]{Corollary}
\newtheorem{conjecture}[theorem]{Conjecture}
\newtheorem{certificate}[theorem]{Certificate}
\newtheorem{definitionx}[theorem]{Definition}
\newtheorem{remark}[theorem]{Remark}
\DeclareMathOperator{\den}{den}
\newcommand{\Zt}{\mathbb Z_2}
\newcommand{\Qt}{\mathbb Q_2}
\newcommand{\ma}{\mu_\alpha}
\newcommand{\mb}{\mu_\beta}
\newcommand{\Nn}{\mathcal N}
\newcommand{\doi}[1]{\href{https://doi.org/#1}{\nolinkurl{doi:#1}}}
\begin{document}
\title[The dyadic denominator law for the phase constants of the Ja]{The dyadic denominator law for the phase constants of the Jacobi zeros}
\author{Iván Area}
\date{}
\maketitle
\noindent\emph{Source and licence.} The dyadic denominator law for the phase constants of the Jacobi zeros. Version arXiv:2608.23006v1 (CC BY 4.0). This material is distributed under the Creative Commons Attribution 4.0 International licence, http://creativecommons.org/licenses/by/4.0/, as recorded in the arXiv OAI licence metadata for this exact version and shown on the abstract page https://arxiv.org/abs/2608.23006v1. Full rights notice: licenses/arxiv-2608.23006-CC-BY-4.0.txt.\par\smallskip
\noindent\textit{Editorial scope.} Counted unit: the original \texttt{theorem} environment \texttt{thm:transferintro} at source lines 142--164 together with its complete proof at lines 611--626; the delivered document keeps source lines 55--755 so that every referenced label and every citation resolves inside it. Native editable source is the paper's own concatenated TeX; the AMS wrapper is editorial formatting only. No publisher class, upstream script or unrelated artwork is redistributed.
\section{Introduction}

We follow the standard normalisation of Jacobi polynomials and Bernoulli numbers in~\cite{Szego,DLMF}.  Let $P_n^{(\alpha,\beta)}$ be the Jacobi polynomials, $\alpha,\beta>-1$, and write the zeros as $x_{n,k}=\cos\theta_{n,k}$ with $0<\theta_{n,1}<\dots<\theta_{n,n}<\pi$. We shall denote
\[
  \Nn=n+\frac{\alpha+\beta+1}2,\qquad z=\cot\frac\theta2,\qquad   A=\alpha^2,\quad B=\beta^2 .
\]
Transforming the Jacobi equation to Liouville normal form and solving the Kummer phase equation recursively in $\Nn^{-2}$ produces a phase
\begin{equation}\label{eq:phase}
  \Xi(\theta)=\Nn\theta-\Bigl(\frac\alpha2+\frac14\Bigr)\pi   +\sum_{r\ge1}\Nn^{-(2r-1)}\Psi_r(z),
\end{equation}
whose zeros are the solutions of $\Xi(\theta_{n,k})=(k-\frac12)\pi$.  (We reserve $\Phi$ for the Borel transform of Section~\ref{sec:normalform}.)  Each integration in the recursive solution introduces a free additive constant; fixing them by the intrinsic rule ``take the odd Laurent primitive in $z$'' gives densities $\Psi_r$ which are Laurent polynomials in $z$ with coefficients in $\mathbb Q[A,B]$, and the additive constants are exactly
\[
  \kappa_r=\Psi_r(1),
\]
the value at the midpoint $\theta=\pi/2$.  They are antisymmetric
\[
\kappa_r(B,A)=-\kappa_r(A,B), 
\]
so $(A-B)\mid\kappa_r$; in particular they vanish identically in the ultraspherical case $\alpha=\beta$, which is why they are invisible for Legendre and Gegenbauer polynomials.

The first three are
\begin{align*}
  \kappa_1&=\frac{A-B}4,\\
  \kappa_2&=\frac{-7(A-B)+2(A^2-B^2)}{96},\\
  \kappa_3&=\frac{781(A-B)}{7680}-\frac{19(A^2-B^2)}{384}+\cdots
\end{align*}
and their denominators are
\[
  4=2^2,\qquad 96=2^5\!\cdot\!3,\qquad 7680=2^9\!\cdot\!15,\qquad   86016=2^{12}\!\cdot\!21,\qquad\dots
\]
The pattern in the two factors is the subject of this paper.

\subsection{Results}

Write $s_2(m)$ for the binary digit sum, $\nu_2$ for the $2$-adic valuation,
and
\[
  H_r=r+\nu_2\bigl((r-1)!\bigr),\quad
  E_r=(2r-1)+H_r=3r-1+\nu_2\bigl((r-1)!\bigr),\quad
  \omega(s)=\Bigl\lceil\frac s2\Bigr\rceil+\nu_2(s).
\]

\begin{theorem}[odd part]\label{thm:odd}
The odd part of $\den\kappa_r$ divides $\operatorname{lcm}(1,3,\dots,2r-1)$.
\end{theorem}

\begin{theorem}[coarse integrality]\label{thm:coarseintro}
$2^{E_r}\kappa_r\in\Zt[\ma,\mb]$, where $\ma=A-\frac14$, $\mb=B-\frac14$.
\end{theorem}

\begin{theorem}[the $2$-adic vanishing]\label{thm:vanishintro} We have
\[
\displaystyle\sum_{k\ge0}\frac{k!}{(2k+1)!!} =\sum_{k\ge0}\frac{2^k}{(2k+1)\binom{2k}{k}}=0 \quad \text{in} \quad \Qt.
\]
\end{theorem}

\begin{theorem}[binomial-moment valuation law]\label{thm:extremalintro}
For every $m\ge0$,
\[
  \nu_2\Bigl(\sum_{j=0}^{m}\binom mj\frac1{2j+1}\Bigr)=m+\nu_2(m+1).
\]
\end{theorem}

\begin{corollary}[extremal unit, conditional on the top layer]\label{cor:extremalintro}
Combined with the top-layer identification
\[
[\ma^r]\kappa_r=\gamma_r\sum_{j=0}^{r-1}\binom{r-1}j\frac1{2j+1},
 \quad \gamma_r=2\binom{2r}r/\bigl((2r-1)16^r\bigr)
 \]
(Proposition~\ref{prop:appell}, established in exact arithmetic for
$r\le20$),
\[
\nu_2\bigl([\ma^r]\kappa_r\bigr)=-2r-\nu_2\bigl((r-1)!\bigr):
\]
the extremal coefficient of $\kappa_r$ in the Bessel basis is a $2$-adic unit after multiplication by $2^{2r+\nu_2((r-1)!)}$, that is, $r-1$ powers of $2$ short of the coarse bound 
\[
E_r=3r-1+\nu_2((r-1)!)
\]
of Theorem~\ref{thm:coarseintro}.
\end{corollary}

The sharp form of the law concerns the \emph{physical} basis $(A,B)$ rather than the Bessel basis $(\ma,\mb)$.  Let
\[
  \ell_r=[A^1B^0]\kappa_r
\]
be the Legendre-tangent coefficients.  The sharp law is $\nu_2(\ell_r)=-E_r$; within this paper it enters as a conjecture, and it is proved in the companion paper~\cite{Ultraspherical}.  Our main structural result replaces it by a single integrality statement for one even power series.


\begin{theorem}[normal form and transfer]\label{thm:transferintro}
Let
\[
g_n=\partial_AL_n(x_0,y_0;A-\tfrac14,B-\tfrac14)\big|_{A=B=0}
\]
be the coefficients of the global Legendre tangent,
\[
\Phi(t)=\sum_{n\ge1}g_nt^n/(n-1)!
\]
its Borel transform, and define
\[
  W(t)=\frac{\sinh 2t}{t^2}\,\operatorname{Im}\Phi(t),\qquad   w_m=[t^{2m}]W(t),
\]
where $\operatorname{Im}\Phi$ is the coefficientwise imaginary part of
$\Phi$.  Then, with $\tau(m)=-2\nu_2\bigl((2m)!\bigr)$,
\[
  \nu_2(\ell_r)\ge-E_r\ \ \forall r   \quad\Longleftrightarrow\quad   \nu_2(w_m)\ge\tau(m)\ \ \forall m ,
\]
and equality transfers index by index.  Equivalently, the sharp law is
\[
  \bigl((2m)!\bigr)^2w_m\in\Zt^\times\qquad\text{for every }m\ge0 .
\]
\end{theorem}


The passage $\Phi\mapsto W$ is an infinite resummation which is nevertheless exact on valuations; this is the first transformation in this problem with that property, and Section~\ref{sec:normalform} explains why every earlier, local, attempt failed.

\begin{conjecture}[W]\label{conj:W}
$\bigl((2m)!\bigr)^2w_m$ is a $2$-adic unit for every $m\ge0$.
\end{conjecture}

\medskip\noindent\emph{Note added.}  Conjecture~\ref{conj:W} (and with it, by Theorem~\ref{thm:transferintro}, the sharp law $\nu_2(\ell_r)=-E_r$ for every $r$) is \emph{proved} in~\cite{Ultraspherical}: the transverse tangent is computed there in closed form (a trigamma kernel for its real part, an explicit Genocchi--Bernoulli central-difference quadrature for its imaginary part, proved in all orders with complete sectorial estimates), and an elementary $2$-adic dominance argument  (a Genocchi-form exponent whose unique $2$-adically dominant partition is the all-ones one) gives $\nu_2(w_m)=-2\nu_2((2m)!)$ for every $m$.  The statements and proofs of the present paper are unchanged and self-contained; where the text below refers to the conjecture as open, this note supersedes it.

The top-layer identity used to convert the binomial-moment valuation into the extremal-unit statement is established by exact arithmetic for $r\le20$. Beyond that range, the extremal-unit conclusion is conditional on that identity, whereas the binomial-moment valuation law itself holds in every order.

Conjecture~\ref{conj:W} is proved in~\cite{Ultraspherical}; independently of that proof, it is verified here for $m\le46$, equivalently $\nu_2(\ell_r)=-E_r$ for $r\le47$, and $\nu_2(\ell_r)=-E_r$ is checked directly for $r\le48$ (Section~\ref{sec:evidence}).  We also record a negative result which we believe is the most informative experimental fact about $W$: no linear ODE with polynomial coefficients, no $P$-recursion and no algebraic equation of size up to $35$ annihilates $W$ or its factorially normalised transform.  A Dwork--Frobenius integrality argument~\cite{Dwork,Nemoto}  therefore cannot be started by guessing an operator.

\subsection{Reproducibility}

All statements below marked \emph{verified} or \emph{certified} are exact computations in rational or Gaussian-rational arithmetic; none is floating point.  The code is reproduced in Section~\ref{sec:code} and supplied as a package.

\section{Setting}\label{sec:setting}

\subsection{The inverse-factorial solution and the Horn recurrence}

Let 
\[
t=e^{i\theta}, \quad w=2/(t-1), \quad \widetilde w=2/(t+1), 
\]
and let
\[
  a_k(\nu)=\frac{(\frac12+\nu)_k(\frac12-\nu)_k}{(-2)^kk!} =\prod_{j=0}^{k-1}\frac{4\nu^2-(2j+1)^2}{8^kk!}
\]
be the Hankel coefficients.  The Jacobi asymptotics at large degree is governed by the inverse-factorial series
\begin{equation}\label{eq:invfact}
  F=\sum_{n\ge0}\frac{G_n}{(Z-1)(Z-2)\cdots(Z-n)},\qquad   G_n=\sum_{l=0}^na_l(\alpha)a_{n-l}(\beta)(-w)^l\widetilde w^{\,n-l},
\end{equation}
with $Z=2\Nn$.  Introduce the Horn variables
\[
  x=-\frac1{t-1}=-\frac w2,\qquad y=\frac1{t+1}=\frac{\widetilde w}2 ;
\]
at the midpoint $\theta=\pi/2$ one has $t=i$ and
\[
  x_0=\frac{1+i}2,\qquad y_0=\frac{1-i}2 ,
\]
so that $-w=1+i$ and $\widetilde w=1-i$ there.  Writing $L=\log F$ and $L=\sum_{n\ge1}L_nZ^{-n}$, we take as the working definition of the phase constants
\begin{equation}\label{eq:kappafromL}
  \kappa_r=2^{-(2r-1)}\operatorname{Im}L_{2r-1}(x_0,y_0).
\end{equation}

\begin{remark}[amplitude--phase identification]\label{rem:identification}
Formula \eqref{eq:kappafromL} is the amplitude--phase dictionary of the Liouville normal form: the midpoint value of the odd Laurent primitive of the $r$-th phase density is carried by the imaginary part of the resummed amplitude logarithm.  In this paper \eqref{eq:kappafromL} is taken as the definition of $\kappa_r$; its identification with the constants $\Psi_r(1)$ of \eqref{eq:phase} is verified in exact rational arithmetic for $r\le8$ against an independent Stirling--Newton implementation of the phase recursion (both engines are included in the package), and the values $\kappa_1,\kappa_2,\kappa_3$ displayed in the introduction agree with both.
\end{remark}

The resummation of \eqref{eq:invfact} is a Kamp\'e de F\'eriet double series in $(x,y)$; it satisfies a Horn system, and its logarithm satisfies a Riccati form which is \emph{division-free}.  Put 
\[
U=\Theta_xL, \quad V=\Theta_yL, \quad T=U+V,
\]
with $\Theta_x=x\partial_x$, $\Theta_y=y\partial_y$.

\begin{theorem}[integral logarithmic recurrence]\label{thm:hornrec}
Let
\[
U=\sum_{n\ge1}U_n\zeta^n, \quad V=\sum_{n\ge1}V_n\zeta^n, \quad \zeta=Z^{-1}.
\]
Then, for $n\ge0$,
\begin{equation}\label{eq:hornrecc}
  U_{n+1}=\Theta_xT_n+\sum_{a+b=n}U_aT_b -x\Bigl(\Theta_xU_n+\sum_{a+b=n}U_aU_b+U_n-\ma\delta_{n,0}\Bigr),
\end{equation}
and symmetrically for $V$ with $x\leftrightarrow y$, $\ma\leftrightarrow\mb$
(the sums run over $a,b\ge1$).  Consequently
\[
U_n,V_n\in\mathbb Z[\ma,\mb][x,y]
\]
for every $n$, and the support obeys $1\le a+b\le c+d\le n$ for every monomial $x^cy^d\ma^a\mb^b$ occurring in $U_n$ or $V_n$.
\end{theorem}

\begin{proof}
Since $1/\bigl((Z-1)\cdots(Z-n)\bigr)=(-1)^n/(1-Z)_n$ and
\[
  a_\mu(\alpha)\,(2x)^\mu(-1)^\mu
  =\frac{(\tfrac12-\alpha)_\mu(\tfrac12+\alpha)_\mu}{\mu!}\,x^\mu ,
\]
term by term in \eqref{eq:invfact}, the series $F$ is the double hypergeometric series
\begin{equation}\label{eq:KdF}
  F=\sum_{\mu,\nu\ge0}
  \frac{(\tfrac12-\alpha)_\mu(\tfrac12+\alpha)_\mu
        (\tfrac12-\beta)_\nu(\tfrac12+\beta)_\nu}
       {(1-Z)_{\mu+\nu}\,\mu!\,\nu!}\;x^\mu y^\nu ,
\end{equation}
an identity of formal series whose coefficients are rational functions of $Z$.  Write $F_{\mu\nu}$ for the coefficient of $x^\mu y^\nu$ in
\eqref{eq:KdF}.  From
\[
  \frac{F_{\mu+1,\nu}}{F_{\mu\nu}}
  =\frac{(\tfrac12-\alpha+\mu)(\tfrac12+\alpha+\mu)}
        {(1-Z+\mu+\nu)(\mu+1)}
  =\frac{(\mu+\tfrac12)^2-A}{(1-Z+\mu+\nu)(\mu+1)}
\]
one reads off, comparing coefficients of $x^{\mu+1}y^\nu$, the partial differential equation
\begin{equation}\label{eq:hornpde}
  \Theta_x\bigl(\Theta_x+\Theta_y-Z\bigr)F
  =x\Bigl(\bigl(\Theta_x+\tfrac12\bigr)^2-A\Bigr)F ,
\end{equation}
and symmetrically in $y$ with $B$.  Substituting $F=e^L$ and dividing by $F$ (legitimate, since $F=1+O(\zeta)$ as a formal series), one has $(\Theta_x+\Theta_y)F=TF$, $\Theta_x(TF)=(\Theta_xT+TU)F$, $\Theta_x(ZF)=ZUF$ and $(\Theta_x+\tfrac12)^2F=(\Theta_xU+U^2+U+\tfrac14)F$, so \eqref{eq:hornpde} becomes the division-free Riccati form \begin{equation}\label{eq:riccati}
  ZU=\Theta_xT+UT-x\bigl(\Theta_xU+U^2+U-\ma\bigr),
\end{equation}
using $A-\tfrac14=\ma$.  Comparing coefficients of $\zeta^n$ (the left side has $\zeta^n$-coefficient $U_{n+1}$) gives \eqref{eq:hornrecc}, with $U_1=x\ma$, $V_1=y\mb$ at $n=0$.  Integrality and support follow by induction: they hold for $U_1,V_1$; every operation on the right-hand side of \eqref{eq:hornrecc}  ($\Theta_x$, multiplication by $x$, integer sums and products) preserves $\mathbb Z[\ma,\mb][x,y]$; in a product the $\mu$-degrees and the radial degrees both add, so $a+b\le c+d$ is preserved; $\Theta_x$ preserves both degrees, the outer multiplication by $x$ raises the radial degree by one, and every term of $U_{n+1}$ has radial degree at most $n+1$ and $\mu$-degree at least $1$.
\end{proof}

No division occurs anywhere in the recurrence; this is what makes the denominators of $\kappa_r$ accessible.  Since $L_n=T_n/(c+d)$ coefficientwise, all denominators are produced by the single division by the radial degree $c+d$ and by the evaluation at $(x_0,y_0)$.

\subsection{The midpoint cost}

Let $K_{c,d}=x_0^cy_0^d/(c+d)$ be the weight attached to a monomial of radial degree $s=c+d$.

\begin{lemma}[exact midpoint cost]\label{lem:cost}
It holds
\[
\nu_2\bigl(\operatorname{Im}K_{c,d}\bigr)=-\omega(s)
\]
whenever $\operatorname{Im}(x_0^cy_0^d)\ne0$, where $\omega(s)=\lceil s/2\rceil+\nu_2(s)$.
\end{lemma}

\begin{proof}
Write $k=\min(c,d)$ and $m=|c-d|$, so $s=2k+m$.  Since $(1+i)(1-i)=2$,
\[
  x_0^cy_0^d=2^{-s}(1+i)^c(1-i)^d=2^{-s+k}\,\zeta^{m},  \qquad \zeta=1\pm i ,
\]
and $\zeta^m=2^{m/2}e^{\pm im\pi/4}$ has $\operatorname{Im}\zeta^m=\pm2^{m/2}\sin(m\pi/4)$.  Hence
\[
\nu_2(\operatorname{Im}\zeta^m) = \begin{cases} \frac{m-1}2,  & m \text{ odd}, \\  \frac m2, & m\equiv2\pmod4, \end{cases}
\]
  and $\operatorname{Im}\zeta^m=0$ for $m\equiv0\pmod4$.  In the first case $s$ is odd and
\[
  \nu_2\bigl(\operatorname{Im}x_0^cy_0^d\bigr)  =-s+k+\frac{m-1}2=-\frac{s+1}2=-\Bigl\lceil\frac s2\Bigr\rceil ,
\]
while $\nu_2(s)=0$; in the second case $s$ is even and
\[
\nu_2(\operatorname{Im}x_0^cy_0^d)=-s+k+\frac m2=-\frac s2 =-\lceil s/2\rceil.
\]
Dividing by $s$ subtracts a further $\nu_2(s)$ in both cases, which is the assertion.
\end{proof}

\begin{lemma}\label{lem:omega}
$\max_{1\le s\le 2r-1}\omega(s)\le H_r=r+\nu_2\bigl((r-1)!\bigr)$.
\end{lemma}

\begin{proof}
By Legendre's formula
\[
H_r=r+\bigl(r-1-s_2(r-1)\bigr)=2r-1-s_2(r-1).
\]
If $s$ is odd then $\nu_2(s)=0$ and $\omega(s)=\lceil s/2\rceil\le r$, while $H_r\ge r$ because $s_2(r-1)\le r-1$.  If $s$ is even then $s\le2r-2$, so
$\omega(s)=\frac s2+\nu_2(s)\le(r-1)+\log_2 s\le(r-1)+\log_2(2r-2)$, and it suffices that
\[
  \log_2(2r-2)+s_2(r-1)\le r ,
  \qquad\text{i.e.}\qquad \log_2m+s_2(m)\le m\quad(m=r-1).
\]
Since $s_2(m)\le\lfloor\log_2m\rfloor+1$, this holds as soon as $2\log_2m+1\le m$, that is for $m\ge7$; the finitely many cases $r\le7$ are
checked directly ($H_r=1,2,4,5,8,9,11$ against $\max_s\omega(s)=1,2,4,4,7,7,8$).
\end{proof}

Theorem~\ref{thm:coarseintro} follows: by Theorem~\ref{thm:hornrec} the numerators are integers, and by Lemmas~\ref{lem:cost}--\ref{lem:omega} the only denominators are $2^{2r-1}$ from \eqref{eq:kappafromL} and $2^{\omega(s)}\le2^{H_r}$ from the midpoint, whence $2^{(2r-1)+H_r}\kappa_r=2^{E_r}\kappa_r\in\Zt[\ma,\mb]$.

Theorem~\ref{thm:odd} follows from the same source: the only odd primes introduced are those of the radial divisor $s=c+d\le2r-1$, and the imaginary part is nonzero only for odd $s$, so the odd part of the denominator divides $\operatorname{lcm}(1,3,\dots,2r-1)$.

\section{The extremal coefficient and a \texorpdfstring{$2$}{2}-adic vanishing}
\label{sec:extremal}

\subsection{The top layer}

The coefficient of the extremal monomial $\ma^r$ in $\kappa_r$ is governed by an Appell layer.  Write
\[
  I_r=\int_0^1(1+t^2)^{r-1}dt=\sum_{j=0}^{r-1}\frac1{2j+1}\binom{r-1}j .
\]

\begin{proposition}[top layer; verified]\label{prop:appell}
$[\ma^r]\kappa_r=\gamma_rI_r$ and $[\mb^r]\kappa_r=-\gamma_rI_r$, with 
\[
\displaystyle\gamma_r=\frac{2\binom{2r}r}{(2r-1)16^r}.
\]
\end{proposition}

Proposition~\ref{prop:appell} is established by exact computation: for $r\le20$ by Lagrange extraction of the leading coefficient of $\kappa_r(A,\tfrac14)$, a polynomial of degree $r$ in $\ma$ (at $\beta=\frac12$ one has $\mb=0$), from exact evaluations of \eqref{eq:kappafromL} at $r+1$ rational parameter values, and for $r\le8$ for the full degree-$r$ homogeneous layer of $\kappa_r(A,B)$.  It is the only computationally established input in this section; everything that follows is proved.

\begin{lemma}[closed form of the binomial moments]\label{lem:endpoint}
Let 
\[
J_m=I_{m+1}=\sum_{j=0}^{m}\binom mj\frac1{2j+1}. 
\]
Then
\[
(2m+1)J_m=2^m+2mJ_{m-1} \quad \text{for } m\ge1, 
\]
and
\[
  J_m=\frac{2^mV_m}{(2m+1)\binom{2m}m},\qquad V_m=\sum_{k=0}^m2^{m-k}\binom{2k}k\in\mathbb Z .
\]
In particular
\[
\nu_2(J_m)=m+\nu_2(V_m)-s_2(m).
\]
\end{lemma}

\begin{proof}
Integration by parts
\[
\frac{d}{dt}\bigl[t(1+t^2)^m\bigr]=(2m+1)(1+t^2)^m-2m(1+t^2)^{m-1}, 
\]
so integrating over $[0,1]$ gives $(2m+1)J_m=2^m+2mJ_{m-1}$.  Setting 
\[
c_m=m!/(2m+1)!!=2^m/\bigl((2m+1)\binom{2m}m\bigr), 
\]
which satisfies $(2m+1)c_m=mc_{m-1}$, the ratio $R_m=J_m/c_m$ obeys $R_m=2R_{m-1}+\binom{2m}m$ with $R_0=1$, whence $R_m=\sum_{k\le m}2^{m-k}\binom{2k}k=V_m$ and $J_m=c_mV_m$.  Kummer's theorem~\cite{Kummer,Mihet,Koblitz} gives $\nu_2\binom{2m}m=s_2(m)$, hence
the valuation formula.
\end{proof}

The target of this section is the valuation law
\begin{equation}\label{eq:lemmaV}
  \nu_2(J_m)=m+\nu_2(m+1),\qquad\text{equivalently}\qquad   \nu_2(V_m)=s_2(m)+\nu_2(m+1),
\end{equation}
which is Theorem~\ref{thm:extremalintro}.

\subsection{Reduction to one series}

Let
\[
  c_k=\frac{k!}{(2k+1)!!}=\frac{2^k}{(2k+1)\binom{2k}k},\qquad
  T_j=\sum_{k\ge j}c_k\quad(j\ge0).
\]

\begin{lemma}[convergence and a floor]\label{lem:conv}
$\nu_2(c_k)=k-s_2(k)=\nu_2(k!)$, which is non-decreasing in $k$ and tends to infinity.  Hence every tail $T_j$ converges in $\Qt$, all $c_k$ have odd
denominators, and
\[
  \nu_2(T_j)\ \ge\ \min_{k\ge j}\nu_2(c_k)\ =\ j-s_2(j).
\]
\end{lemma}

\begin{proof}
Kummer's theorem~\cite{Kummer,Mihet,Koblitz} gives 
\[
\nu_2\binom{2k}k=s_2(k), 
\]
and $2k+1$ is odd, so $\nu_2(c_k)=k-s_2(k)$, which is Legendre's formula for $\nu_2(k!)$ and is therefore non-decreasing.
\end{proof}

\begin{theorem}[reduction]\label{thm:Freduction}
The valuation law \eqref{eq:lemmaV} holds for every $m$ if and only if
\[
  T_0=\sum_{k\ge0}\frac{k!}{(2k+1)!!}=0\quad\text{in }\Qt .
\]
\end{theorem}

\begin{proof}
Since 
\[
(\tfrac12)_k/(\tfrac32)_k=1/(2k+1), 
\]
one has
\[
J_m={}_2F_1(-m,\tfrac12;\tfrac32;-1), 
\]
and Pfaff's transformation~\cite[Ch.~2]{AndrewsAskeyRoy} with $z=-1$, $z/(z-1)=\tfrac12$ gives $J_m=2^m\,{}_2F_1(-m,1;\tfrac32;\tfrac12)$, that is, with $B_m=J_m/2^m$,
\[
  B_m={}_2F_1\bigl(-m,1;\tfrac32;\tfrac12\bigr) =\sum_{k\ge0}(-1)^k\binom{m}{k}c_k ,
\]
because $(-m)_k(1)_k2^{-k}/\bigl((\tfrac32)_kk!\bigr)=(-1)^k\binom mkc_k$. The law \eqref{eq:lemmaV} is $\nu_2(B_m)=\nu_2(m+1)$.

Write $n=m+1$ and $\binom mk=\binom{n-1}k=\sum_{j\le k}(-1)^{k-j}\binom nj$; exchanging the order of summation (legitimate in $\Qt$ since $c_k\to0$),
\[
  B_m=\sum_{j\ge0}(-1)^j\binom njT_j =T_0+n\sum_{j\ge1}(-1)^j\,\frac{T_j}{j}\binom{n-1}{j-1},
\]
using $\binom nj=\frac nj\binom{n-1}{j-1}$ for $j\ge1$.

Assume $T_0=0$.  By Lemma~\ref{lem:conv}
\[
\nu_2(T_j)\ge j-s_2(j)=\nu_2(j!),
\]
and 
\[
\nu_2(j!)=\nu_2(j)+\nu_2\bigl((j-1)!\bigr)\ge\nu_2(j)+1
\] 
for $j\ge3$; hence $T_j/j\in2\Zt$ for $j\ge3$.  For $j=2$, the exact value $T_2=T_0-c_0-c_1=-\tfrac43$ gives $T_2/2\in2\Zt$; and $T_1=T_0-c_0=-1$, so the $j=1$ term equals $c_0\binom{n-1}0=1$.  The bracketed sum is therefore a $2$-adic unit, and $\nu_2(B_m)=\nu_2(n)=\nu_2(m+1)$, which is \eqref{eq:lemmaV}.

Conversely, assume \eqref{eq:lemmaV}.  Since $c_k\to0$ in $\Zt$, the map $m\mapsto B_m$ is the restriction to non-negative integers of a Mahler series, hence extends continuously to $\Zt$~\cite[Ch.~4]{Koblitz}.  As $2^k-1\to-1$ in $\Zt$ and $\nu_2(B_{2^k-1})=\nu_2(2^k)=k\to\infty$,
continuity gives $B_{-1}=\sum_{k\ge0}(-1)^k\binom{-1}kc_k=\sum_{k\ge0}c_k=T_0=0$.
\end{proof}

\subsection{The vanishing}

For $i\ge0$ let
\[
  S_i=\sum_{k\ge i}c_k\binom ki
\]
be the \emph{binomial} moments of the series, convergent in $\Qt$ by Lemma~\ref{lem:conv}, with $\nu_2(S_i)\ge i-s_2(i)$ and $S_0=T_0$.  (The power moments would not satisfy the clean recurrence below.)

\begin{theorem}\label{thm:vanish}
$S_0=0$ in $\Qt$.
\end{theorem}

\begin{proof}
The contiguity $(2k+1)c_k=kc_{k-1}$, valid for $k\ge1$, gives
\[
  \sum_kc_k(2k+1)\binom ki=\sum_kkc_{k-1}\binom ki =\sum_k(k+1)c_k\binom{k+1}i ,
\]
that is
\[
\sum_kc_k\bigl[(2k+1)\binom ki-(k+1)\binom{k+1}i\bigr]=0.
\]
Now
\[
\binom{k+1}i=\binom ki+\binom k{i-1}, \quad (k-i)\binom ki=(i+1)\binom k{i+1}, \quad (k-i+1)\binom k{i-1}=i\binom ki,
\]
so the bracket equals $(i+1)\binom k{i+1}-i\binom k{i-1}$ and therefore
\begin{equation}\label{eq:contig}
  (i+1)S_{i+1}=iS_{i-1}\qquad(i\ge1).
\end{equation}
Taking $i=2j-1$ in \eqref{eq:contig} gives $2jS_{2j}=(2j-1)S_{2j-2}$, so by induction
\[
  S_{2j}=S_0\prod_{t=1}^j\frac{2t-1}{2t}=\frac{S_0}{4^j}\binom{2j}j,   \qquad\text{whence}\qquad   \nu_2(S_{2j})=\nu_2(S_0)+s_2(j)-2j .
\]
On the other hand Lemma~\ref{lem:conv} at $i=2j$ gives
\[
\nu_2(S_{2j})\ge 2j-s_2(2j)=2j-s_2(j).
\]
Combining,
\[
  \nu_2(S_0)\ \ge\ 4j-2s_2(j)\qquad\text{for every }j\ge1 ,
\]
and the right-hand side is unbounded.  Hence $S_0=0$.
\end{proof}

\begin{corollary}\label{cor:F}
The valuation law \eqref{eq:lemmaV} holds for every $m$, which is Theorem~\ref{thm:extremalintro}; and $\nu_2(I_r)=(r-1)+\nu_2(r)$ for every $r\ge1$.
\end{corollary}

\begin{proof}
Theorems~\ref{thm:Freduction} and~\ref{thm:vanish}, then $I_r=J_{r-1}$.
\end{proof}

\begin{corollary}[extremal unit, conditional on Proposition~\ref{prop:appell}]\label{cor:extremal}
With the top-layer identification of Proposition~\ref{prop:appell},
\[
\nu_2\bigl([\ma^r]\kappa_r\bigr)=\nu_2(\gamma_r)+\nu_2(I_r)=-2r-\nu_2\bigl((r-1)!\bigr),
\]
which is Corollary~\ref{cor:extremalintro}; the identity of exponents uses
\[
\nu_2(\gamma_r)=1+s_2(r)-4r, 
\]
Corollary~\ref{cor:F}, and
\[
\nu_2(r)=1-s_2(r)+s_2(r-1)
\]
 together with Legendre's formula $\nu_2((r-1)!)=r-1-s_2(r-1)$.
\end{corollary}

\begin{remark}[the Bessel-basis denominator gap]\label{rem:besselgap}
Granting Proposition~\ref{prop:appell}, the Bessel-basis denominator of $\kappa_r$ satisfies
\[
  2r+\nu_2\bigl((r-1)!\bigr)\ \le\
  \nu_2\bigl(\den\nolimits_\mu\kappa_r\bigr)\ \le\
  3r-1+\nu_2\bigl((r-1)!\bigr),
\]
the upper bound by Theorem~\ref{thm:coarseintro} and the lower bound by Corollary~\ref{cor:extremal}; the two differ by exactly $r-1$.  Closing this gap is open.
\end{remark}

\begin{remark}
For related work on $p$-adic valuations and approximations of binomial sums, see~\cite{AidagulovAlekseyev,ZhangYuan}.  Theorem~\ref{thm:vanish} looks like a curiosity but it is the arithmetic heart of the extremal computation, and we have not found it in the literature.  Its real-analytic counterpart is the classical
\[
\sum_k2^k/\bigl((2k+1)\binom{2k}k\bigr)=\pi/2;
\]
the $2$-adic sum of the same series vanishes.  The partial sums $\sum_{k<K}c_k$ have $\nu_2$ equal to $9$, $20$, $48$, $99$, $200$, $401$, $802$ for $K=10$, $20$, $50$, $100$, $200$, $400$, $800$, just above the floor $\nu_2(c_K)=K-s_2(K)$ of Lemma~\ref{lem:conv}: the tail cancels the head exactly.
\end{remark}

\section{The normal form}\label{sec:normalform}

\subsection{Why localisation fails}

Before stating the transfer we record, because it shapes everything that follows, that the sharp law resists every local treatment.  All of the
following were tried and all fail in the same way, namely $\nu_2\bigl(\sum_sX_s\bigr)>\min_s\nu_2(X_s)$: radial shells; Cartier tails; local Ward currents built from $13$, $33$ and $75$ generators on a nonlinear Horn background with first and second jets and background factors (the
deficiency is exactly one in every case, and adjoining the nonlocal column closes it immediately); residue moments modulo $8$; and boundary jets.  The quantity must be summed before its valuation is taken.  The normal form below is the first construction that performs the whole summation and still controls the valuation.

\subsection{The global tangent and its Borel transform}

Let
\[
  g_n=\partial_AL_n\bigl(x_0,y_0;A-\tfrac14,B-\tfrac14\bigr)\Big|_{A=B=0},
  \qquad
  G(\zeta)=\sum_{n\ge1}g_n\zeta^n=\frac{\partial_AF}F\bigg|_{A=B=0},
\]
so that, by \eqref{eq:kappafromL},
\begin{equation}\label{eq:ltangent}
  \ell_r=2^{-(2r-1)}\operatorname{Im}g_{2r-1},\qquad   K(z)=\sum_{r\ge1}\ell_rz^{2r-1}      =\operatorname{Im}\frac{G(z/2)-G(-z/2)}2 .
\end{equation}
Because $g_n$ grows like $(n-1)!$, the natural analytic object is the Borel transform 
\[
\Phi(t)=\sum_{n\ge1}g_nt^n/(n-1)!, 
\]
with coefficients in $\mathbb Q(i)$; we write 
\[
\operatorname{Im}\Phi(t)=\sum_{n\ge1}(\operatorname{Im}g_n)t^n/(n-1)!
\]
for its coefficientwise imaginary part.

\subsection{The transfer theorem}

\begin{definitionx}\label{def:W}
Define
\begin{equation}\label{eq:normalform}
  W(t)=\frac{\sinh2t}{t^2}\,\operatorname{Im}\Phi(t),
  \qquad
  w_m=[t^{2m}]\,W(t),
\end{equation}
so that
\[
  \sum_{m\ge0}w_mt^{2m}
  =1+\frac{t^2}{12}+\frac{t^4}{64}-\frac{143\,t^6}{80640}       +\frac{21787\,t^8}{46448640}-\cdots
\]
\end{definitionx}

Only the even part of $W$ enters the results below; the odd-power coefficients of $W$ are generated by the even-index $\operatorname{Im}g_{2k}$ and vanish in the whole computed range (Remark~\ref{rem:identityR}).

\begin{lemma}[the two weights]\label{lem:weights}
Put
\[
D(t)=\frac{2t}{\sinh2t}=\sum_kd_kt^{2k},\qquad S(t)=\frac{\sinh2t}{2t}=\sum_ke_kt^{2k},
\]
so that $DS=1$.  Then $d_0=e_0=1$ and, for $k\ge1$,
\[
  \nu_2(d_k)=\nu_2(e_k)=s_2(k).
\]
\end{lemma}

\begin{proof}
Legendre's formula and $e_k=2^{2k}/(2k+1)!$  give 
\[
\nu_2(e_k)=2k-\bigl(2k+1-s_2(2k+1)\bigr)=s_2(k), 
\]
using $s_2(2k+1)=s_2(k)+1$.  For $D$, the classical expansion 
\[
x/\sinh x=\sum_k(2-2^{2k})B_{2k}x^{2k}/(2k)!, 
\]
obtained from the Laurent series for $\csc$ in~\cite[Eq.~4.19.4]{DLMF} by $z\mapsto ix$ (cf.~\cite[Eqs.~4.28.8--4.28.13 and 24.2.1]{DLMF}), at $x=2t$ gives $d_k=(2-2^{2k})B_{2k}2^{2k}/(2k)!$; for $k\ge1$, $\nu_2(2-2^{2k})=1$ and $\nu_2(B_{2k})=-1$ by the von Staudt--Clausen theorem~\cite{Clausen,Rado,Carlitz}, while 
\[
\nu_2((2k)!)=2k-s_2(k), 
\]
so $\nu_2(d_k)=s_2(k)$.
\end{proof}


\begin{proof}[Proof of Theorem~\ref{thm:transferintro}]
Write 
\[
\tau(m)=-2\bigl(2m-s_2(m)\bigr)=-2\nu_2((2m)!).  
\]
By Definition~\ref{def:W}, $W=\bigl(2\operatorname{Im}\Phi/t\bigr)S$ and hence $\operatorname{Im}\Phi=\frac t2WD$.  Since $S$ and $D$ are even, the even part of $W$ is the convolution of $S$ with the even part of $2\operatorname{Im}\Phi/t$, and conversely; so with
$b_m=[t^{2m}]\bigl(2\operatorname{Im}\Phi/t\bigr) =2\operatorname{Im}g_{2m+1}/(2m)!$ the sequences $(w_m)$ and $(b_m)$ are related by the mutually inverse convolutions with $(e_k)$ and $(d_k)$, and the odd-index data never enter.  By \eqref{eq:ltangent}, $b_m=2^{2m+2}\ell_{m+1}/(2m)!$, whence
\[
  \nu_2(b_m)=2m+2+\nu_2(\ell_{m+1})-\bigl(2m-s_2(m)\bigr)            =2+s_2(m)+\nu_2(\ell_{m+1}),
\]
and $\nu_2(b_m)\ge\tau(m)$ is exactly $\nu_2(\ell_{m+1})\ge-E_{m+1}$, since $-E_{m+1}=-4m-2+s_2(m)$.  It therefore suffices to show that convolution with a series $c$ satisfying $\nu_2(c_0)=0$ and $\nu_2(c_k)\ge s_2(k)$ for $k\ge1$ preserves the family of bounds $\tau$ and preserves equality.  For $1\le k\le m$, subadditivity $s_2(m)\le s_2(m-k)+s_2(k)$ gives
\[
  \bigl[s_2(k)+\tau(m-k)\bigr]-\tau(m)  =s_2(k)+4k+2s_2(m-k)-2s_2(m)\ \ge\ 4k-s_2(k)\ >\ 0,
\]
because $s_2(k)\le\log_2k+1<4k$ for $k\ge1$.  So in $\sum_kc_ka_{m-k}$ every term with $k\ge1$ has valuation strictly larger than $\tau(m)$; hence $\nu_2$ of the sum equals $\nu_2(c_0a_m)=\nu_2(a_m)$ whenever $\nu_2(a_m)=\tau(m)$, and is $\ge\tau(m)$ whenever $\nu_2(a_j)\ge\tau(j)$ for $j\le m$.  Applying this to $c=D$ and to $c=S$, both admissible by Lemma~\ref{lem:weights}, gives the equivalence in both directions.
\end{proof}


\begin{remark}[what has changed]
Three things.  First, the open statement is no longer about a family of Horn coefficients indexed by two degrees and two $\mu$-exponents, but about the coefficients of a \emph{single} even power series in one variable.  Second, the weight has become uniform: $\tau(m)=-2\nu_2((2m)!)$ replaces $E_r=3r-1+\nu_2((r-1)!)$.  Third (and this is the methodological point) the passage $\Phi\mapsto W$ is an infinite resummation which is exact on valuations, because the whole tail is strictly subdominant.
\end{remark}

\begin{remark}[the symmetric half of the tangent; verified]\label{rem:identityR}
In the whole computed range $n\le95$ the tangent satisfies, exactly,
\[
  \operatorname{Re}\Phi(t)=\frac{t^2e^t}{\sinh2t} =\frac{t^2}2\bigl(\operatorname{csch}t+\operatorname{sech}t\bigr),
  \qquad \operatorname{Im}g_{2k}=0 ,
\]
equivalently $\operatorname{Re}g_{2k+1}=(1-2^{2k-1})B_{2k}$ and $\operatorname{Re}g_{2k}=\frac{(2k-1)E_{2k-2}}2$, with $B$ the Bernoulli and $E$ the Euler numbers.  Granting these closed forms, $W$ is even, the definition \eqref{eq:normalform} takes the symmetrical shape
\[
  \frac{\sinh2t}{t^2}\,\Phi(t)=e^t+i\,W(t),
\]
and $\nu_2(\operatorname{Re}g_n)=-1$ for every $n\ge1$: Euler numbers are odd integers, so $\nu_2\bigl((2k-1)E_{2k-2}/2\bigr)=-1$; $1-2^{2k-1}$ is odd, and von Staudt--Clausen gives $\nu_2(B_{2k})=-1$; the case $n=1$ is $\operatorname{Re}g_1=\frac12$.  The symmetric half of the tangent therefore carries no dyadic information beyond one unit at each order: all that is open sits in the antisymmetric half.  These closed forms reflect the ultraspherical diagonal of the amplitude; they are recorded here as verified, and \emph{none of the results of this paper depends on them}.
\end{remark}

\section{Evidence, and one negative certificate}\label{sec:evidence}

\subsection{Verification range}

The tangent coefficients $g_n$ were computed exactly for $n\le95$ by a dual-number specialisation of Theorem~\ref{thm:hornrec} (Section~\ref{sec:code}), in pure integer arithmetic at the implicit scale $4^{-n}$.  Within that range:
\[
  \nu_2(\ell_r)=-E_r\ \text{ for }r\le48,\qquad
  \nu_2(w_m)=-2\nu_2\bigl((2m)!\bigr)\ \text{ for }m\le46 .
\]
Table~\ref{tab:units} lists the odd parts of the normalised units.

\begin{table}[t]
\centering
\small
\begin{tabular}{@{}rrl@{}}
\toprule
$m$ & $\nu_2(w_m)$ & odd part of $\bigl((2m)!\bigr)^2w_m$\\
\midrule
$0$ & $0$   & $1$\\
$1$ & $-2$  & $1/3$\\
$2$ & $-6$  & $9$\\
$3$ & $-8$  & $-6435/7$\\
$4$ & $-14$ & $762545$\\
$5$ & $-16$ & $-18081757575/11$\\
$6$ & $-20$ & $8219311479225$\\
$7$ & $-22$ & $-82171171579421925$\\
$8$ & $-30$ & $1483464568148296358625$\\
$9$ & $-32$ & $-851695035559439777297967375/19$\\
$10$& $-36$ & $2137965664907224677330336155625$\\
\bottomrule
\end{tabular}
\caption{The normalised coefficients of Conjecture~\ref{conj:W}.  The valuations follow the ruler sequence exactly:
$\nu_2(w_m)=-2\nu_2((2m)!)$.}
\label{tab:units}
\end{table}

\subsection{No small annihilating operator}

\begin{certificate}[no small closed equation]\label{cert:noeq}
With $g_n$ known exactly for $n\le95$, hence $w_m$ for $m\le46$, none of the three series
\[
  W(z)=\sum_mw_mz^m,\qquad   \mathcal U(z)=\sum_m\bigl((2m)!\bigr)^2w_mz^m,\qquad   \sum_m\frac{\operatorname{Im}g_{2m+1}}{(2m)!}z^m
\]
satisfies a $P$-recursion $\sum_{i\le R}q_i(m)a_{m+i}=0$, a linear ODE $\sum_{i\le R}p_i(z)f^{(i)}=0$, or an algebraic equation $\sum_{i\le R}q_i(z)f^i=0$, with $\deg\le D$ and $(R+1)(D+1)\le35$, at least six spare equations beyond the number of unknowns, and a reserved block of coefficients on which each candidate operator is retested.  The same programme recovers the expected operators for $\sum z^m/((2m)!)^2$, for the Catalan numbers and for $\sum m!\,z^m$, and correctly finds nothing for an $s_2$-driven control series.
\end{certificate}

\begin{remark}
A negative result of this kind does not exclude a closed equation; it excludes the small ones.  But those are exactly the ones a Dwork--Frobenius argument needs.  The consequence for the programme is that a functional equation for $W$, if there is one, has to come from the Horn structure underlying Theorem~\ref{thm:hornrec} rather than from guessing.
\end{remark}

\subsection{The Frobenius form}

Once $\mathcal U\in\Zt[[z]]$ is known, $u_m\in\Zt^\times$ is equivalent to $u_m\equiv1\pmod 2$, because every unit of $\Zt$ is $\equiv1$.  Hence Conjecture~\ref{conj:W} takes the compact form
\[
  (1-z)\,\mathcal U(z)\equiv1\pmod 2 .
\]
This is a faithful restatement and a convenient normal form; it is not a reduction, since the Cartier descent $u_{2m}\equiv u_m$ together with
$u_{2m+1}\equiv1$ is a bookkeeping split of the same statement.  All the content is in the integrality of $\mathcal U$.


\section{A Mellin split for the zero-balanced route}\label{sec:mellinsplit}
A parallel manuscript~\cite{Ultraspherical} reduces the physical tangent to the symmetric parameter variation of a zero-balanced Gauss function at the prescribed upper boundary value $2+i0$: with $s=(1-Z)/2$, the tangent is obtained from ${}_2F_1^{\uparrow}(s-\delta/2,s+\delta/2;2s;2)$ by two derivatives in $\delta$ at $\delta=0$.  The following elementary split makes the branch contribution explicit and gives a concrete starting point for a $2$-adic coefficient analysis.  It does not, by itself, prove Conjecture~\ref{conj:W}; the proof of the conjecture is given in~\cite{Ultraspherical}, along a different route (the zero-balanced boundary split, a two-component matching, and $2$-adic dominance).

\begin{proposition}[upper-boundary Mellin split]\label{prop:mellinsplit}
Put $a=s-\delta/2$, $b=s+\delta/2$, and assume first that $0<\operatorname{Re}a<1$ and $\operatorname{Re}b>0$, so that Euler's integral and both limiting integrals below converge.  Then
\begin{align*}
 {}_2F_1^{\uparrow}(a,b;a+b;2)
 =\frac{\Gamma(a+b)}{\Gamma(a)\Gamma(b)}\bigg\{&
 2^{-b}\int_0^1 u^{b-1}(1-u/2)^{a-1}(1-u)^{-a}\,du\\
 &+e^{i\pi a}2^{1-a-b}\int_0^1
 u^{-a}(1-u)^{a-1}(1+u)^{b-1}\,du\bigg\}.
\end{align*}
The identity extends meromorphically in $(a,b)$.
\end{proposition}
\begin{proof}
Start with Euler's integral for ${}_2F_1(a,b;a+b;z)$, analytic in $z$ on the upper half-plane, and let $z\to2+i0$; dominated convergence applies in the stated strip.  Split the integral at $t=1/2$.  In the first half put $t=u/2$; in the second put $t=(1+u)/2$.  On the second half $1-(2+i0)t$ approaches the negative axis from below, hence $(1-(2+i0)t)^{-a}=e^{i\pi a}(2t-1)^{-a}$ on the upper-boundary branch. The displayed formula follows.  By Corollary~\ref{cor:kernels} below both integrals are Beta multiples of Gauss functions, hence meromorphic in $(a,b)$, and so is the left-hand boundary value; the identity therefore persists by analytic continuation wherever both sides are finite.
\end{proof}

\begin{corollary}[closed form of the two kernels]\label{cor:kernels}
In the strip of Proposition~\ref{prop:mellinsplit},
\[
 \int_0^1u^{b-1}(1-u/2)^{a-1}(1-u)^{-a}\,du
 =B(b,1-a)\,{}_2F_1\bigl(1-a,b;\,b+1-a;\,\tfrac12\bigr),
\]
\[
 \int_0^1u^{-a}(1-u)^{a-1}(1+u)^{b-1}\,du
 =B(1-a,a)\,{}_2F_1\bigl(1-b,1-a;\,1;\,-1\bigr).
\]
At the diagonal point $\delta=0$, i.e.\ $a=b=s$, both Gauss functions are summable in closed form: the first by Gauss's second summation theorem (there $c=(\alpha+\beta+1)/2$) and the second by Kummer's theorem (there $c=1+\alpha-\beta$), see~\cite[Sec.~3.1]{AndrewsAskeyRoy}:
\[
 {}_2F_1\bigl(1-s,s;1;\tfrac12\bigr)
 =\frac{\sqrt\pi}{\Gamma\bigl(1-\frac s2\bigr)
                    \Gamma\bigl(\frac{s+1}2\bigr)},
 \qquad
 {}_2F_1\bigl(1-s,1-s;1;-1\bigr)
 =\frac{\Gamma\bigl(\frac{3-s}2\bigr)}
        {\Gamma(2-s)\,\Gamma\bigl(\frac{1+s}2\bigr)} .
\]
\end{corollary}

\begin{proof}
Both integrals are Euler representations: $\int_0^1u^{\beta-1}(1-u)^{\gamma-\beta-1}(1-zu)^{-\alpha}du =B(\beta,\gamma-\beta)\,{}_2F_1(\alpha,\beta;\gamma;z)$ with $(\alpha,\beta,\gamma,z)=(1-a,\,b,\,b+1-a,\,\tfrac12)$ for the first and $(1-b,\,1-a,\,1,\,-1)$ for the second.  At $a=b=s$ the first Gauss function has parameters $(1-s,s;1)$ with $1=\bigl((1-s)+s+1\bigr)/2$, so Gauss's second theorem applies and gives $\Gamma(\tfrac12)\Gamma(1)/\bigl(\Gamma(1-\tfrac s2) \Gamma(\tfrac{s+1}2)\bigr)$; the second has parameters $(1-s,1-s;1)$ with $1=1+\alpha-\beta$, so Kummer's theorem applies and gives $\Gamma(1)\Gamma(\tfrac{3-s}2)/\bigl(\Gamma(2-s)\Gamma(\tfrac{1+s}2)\bigr)$.
\end{proof}

\begin{remark}[what this attack achieves and what remains]
Proposition~\ref{prop:mellinsplit} separates the boundary value into a branch-free part and one explicit phase $e^{i\pi a}$, and Corollary~\ref{cor:kernels} shows that at $\delta=0$ both parts collapse to Gamma quotients: the zero-balanced boundary value is an explicit perturbation of two classically summable points.  The transverse tangent (two $\delta$-derivatives at $\delta=0$) therefore involves only first and second parameter derivatives of Gauss functions at a
Gauss-second point and at a Kummer point, objects with classical digamma-series expansions, together with the derivative of the explicit phase.  The unresolved step is arithmetic: to turn the resulting large-$Z$ coefficient extraction into an integral lattice whose reduction modulo $2$ proves $((2m)!)^2w_m\equiv1\pmod2$.  The split also confirms that changing $2+i0$ to $2-i0$ conjugates the response and cannot be ignored.  Analytic background on zero-balanced functions is available in~\cite{AndersonEtAl,SimicVuorinen}; the specialization itself is the one classified by Vid\=unas~\cite{VidunasAppell}, and its derivation on the physical curve is given in~\cite{Ultraspherical}.
\end{remark}

\section{The code}\label{sec:code}


\begin{thebibliography}{99}
\bibitem[AidagulovAlekseyev]{AidagulovAlekseyev} R. R. Aidagulov and M. A. Alekseyev, On $p$-adic approximation of sums of binomial coefficients, \emph{J. Math. Sci.} \textbf{233} (2018), 626--634. \doi{10.1007/s10958-018-3948-0}.
\bibitem[AndersonEtAl]{AndersonEtAl} G. D. Anderson, R. W. Barnard, K. C. Richards, M. K. Vamanamurthy and M. Vuorinen, Inequalities for zero-balanced hypergeometric functions, \emph{Trans. Amer. Math. Soc.} \textbf{347} (1995), 1713--1723. \doi{10.1090/S0002-9947-1995-1264800-3}.
\bibitem[AndrewsAskeyRoy]{AndrewsAskeyRoy} G. E. Andrews, R. Askey and R. Roy, \emph{Special Functions}, Encyclopedia of Mathematics and its Applications, vol.~71, Cambridge University Press, Cambridge, 1999. \doi{10.1017/CBO9781107325937}.
\bibitem[Carlitz]{Carlitz} L. Carlitz, The Staudt--Clausen theorem, \emph{Math. Mag.} \textbf{34} (1961), 131--146.
\bibitem[Clausen]{Clausen} [T. Clausen], Theorem, \emph{Astronomische Nachrichten} \textbf{17} (1840), cols.~351--352. \doi{10.1002/asna.18400172204}. (The note is unsigned in \emph{AN}; the attribution to Clausen is the standard one.)
\bibitem[DLMF]{DLMF} F. W. J. Olver, A. B. Olde Daalhuis, D. W. Lozier, B. I. Schneider, R. F. Boisvert, C. W. Clark, B. R. Miller, B. V. Saunders, H. S. Cohl and M. A. McClain, eds., \emph{NIST Digital Library of Mathematical Functions}, Release 1.2.7 of 2026-06-15, \url{https://dlmf.nist.gov/}.
\bibitem[Dwork]{Dwork} B. Dwork, $p$-adic cycles, \emph{Publ. Math. Inst. Hautes \'Etudes Sci.} \textbf{37} (1969), 27--115. \doi{10.1007/BF02684886}.
\bibitem[Koblitz]{Koblitz} N. Koblitz, \emph{$p$-adic Numbers, $p$-adic Analysis, and Zeta-Functions}, 2nd ed., Graduate Texts in Mathematics, vol.~58, Springer, New York, 1984. \doi{10.1007/978-1-4612-1112-9}.
\bibitem[Kummer]{Kummer} E. E. Kummer, \"Uber die Erg\"anzungss\"atze zu den allgemeinen Reciprocit\"atsgesetzen, \emph{J. Reine Angew. Math.} \textbf{44} (1852), 93--146. \doi{10.1515/crll.1852.44.93}.
\bibitem[Mihet]{Mihet} D. Mihet, Legendre's and Kummer's theorems again, \emph{Resonance} \textbf{15} (2010), 1111--1121. \doi{10.1007/s12045-010-0123-4}.
\bibitem[Nemoto]{Nemoto} Y. Nemoto, Transformation formula of Dwork's $p$-adic hypergeometric function, \emph{Bull. Aust. Math. Soc.} \textbf{113} (2026), 553--564. \doi{10.1017/S0004972725100580}.
\bibitem[Rado]{Rado} R. Rado, A new proof of a theorem of V. Staudt, \emph{J. London Math. Soc.} \textbf{s1-9} (1934), 85--88. \doi{10.1112/jlms/s1-9.2.85}.
\bibitem[SimicVuorinen]{SimicVuorinen} S. Simi\'c and M. Vuorinen, Landen inequalities for zero-balanced hypergeometric functions, \emph{Abstr. Appl. Anal.} \textbf{2012} (2012), Article ID 932061, 11 pp. \doi{10.1155/2012/932061}.
\bibitem[Szego]{Szego} G. Szeg\H{o}, \emph{Orthogonal Polynomials}, 4th ed., American Mathematical Society Colloquium Publications, vol.~23, American Mathematical Society, Providence, RI, 1975. ISBN 978-0-8218-1023-1.
\bibitem[Ultraspherical]{Ultraspherical} I. Area, The ultraspherical background of the Jacobi phase expansion, in closed form, submitted, 2026, % TODO: insert the arXiv number once posted %\href{https://arxiv.org/abs/XXXX.XXXXX}{arXiv:XXXX.XXXXX}.
\bibitem[VidunasAppell]{VidunasAppell} R. Vid\=unas, Specialization of Appell's functions to univariate hypergeometric functions, \emph{J. Math. Anal. Appl.} \textbf{355} (2009), 145--163. \doi{10.1016/j.jmaa.2009.01.047}.
\bibitem[ZhangYuan]{ZhangYuan} Y. Zhang and P. Yuan, The $p$-adic valuations of sums of binomial coefficients, \emph{J. Math.} \textbf{2021} (2021), Article ID 9570350, 12 pp. \doi{10.1155/2021/9570350}.
\end{thebibliography}
\end{document}
